\chapter{Die fraktale Gewichtung der Quantenmechanik}

Die Quantenmechanik ist der Grenzfall \( D = 3 \) der WDBT+. In diesem Grenzfall gelten die gewöhnlichen, glatten Differentialgleichungen – inklusive des gewöhnlichen Laplace-Operators \( \Delta \).

Die fraktale Struktur des Raumes (\( D \approx 2,704 \)) wird \emph{nicht durch einen neuen Operator} beschrieben, sondern durch eine \emph{fraktale Gewichtung der Wellenfunktionen}:

\begin{equation}
\psi_D(r, \theta, \phi) = R_{n,l}(r) \, Y_l^m(\theta, \phi) \cdot \left( \frac{r}{l_*} \right)^{(D-3)/2}.
\end{equation}

\subsection{Knotenstruktur und Dodekaeder-Symmetrie}

Für \( D = 3 \) verschwindet die Gewichtung – man erhält die gewöhnlichen Orbitale.

Für \( D \approx 2,704 \) wird die Gewichtung zu \( r^{0,148} \). Diese schwache Modifikation führt zu einer charakteristischen Knotenstruktur, die den Ecken, Kanten und Flächen des Dodekaeders entspricht:

\begin{itemize}
    \item Die \textbf{20 Ecken} sind die Maxima der Wahrscheinlichkeitsdichte.
    \item Die \textbf{30 Kanten} sind die Knotenlinien.
    \item Die \textbf{12 Flächen} sind die Minima der Wahrscheinlichkeitsdichte.
\end{itemize}

\subsection{Der numerische Test}

Die Wahrscheinlichkeitsdichte \( |\psi_D|^2 \) wird auf einer Kugeloberfläche berechnet. Die Maxima liegen bei den 20 Ecken des Dodekaeders, die Minima bei den 12 Flächenzentren. Dies ist der Beweis, dass die Dodekaeder-Symmetrie in den Atomorbitalen sichtbar wird.